\section{Thin and repair stress spaces}
\label{sec:stress-spaces-short}

Fix the exact background from Section~\ref{short-bg:section}. The goal is
its nonzero periodic energy mode $b$ with exponent $0<\sigma<1/4$.
The planar Gaussian eigenpair supplies transient train amplification
and a signed integrated stress response. The three-dimensional pressure
interaction returns a strain to the young rings and closes the loop.
The singular limit compares these response operators as curvature and
viscosity decrease together; the final characteristic zero belongs to
the exact three-dimensional operator, with its full pressure and repair.

We first establish the causal near inverse and all-age moments, then
identify the returned strain and control its errors, and finally solve
the rank-one scalar return equation. The background and geometry remain
fixed as we vary the real Floquet exponent.

Use the exact background
$U=V+r_b$, where $V=U_{\rm raw}+E_\epsilon\alpha_{\rm bg}$
and $r_b$ is the energy repair. Set $c=\nu^{-2}$, $R=\epsilon^{-1}$,
and retain component order $(r,\phi,z)$. In the cylinder chart
\[
 s=R^2/\rho^2,\qquad q=R\log(\rho/R)+R\tau/2,\qquad
 \eta=Rz/\rho,
\]
put $d\varsigma=ds\,dq/c\,d\eta$ and
\[
 H=L^2(s^{-3/2}d\varsigma),\qquad
 H_T=L^2(s^{-1/2}d\varsigma;\operatorname{Sym}_3).
\]
The natural velocity and tensor lifts are
$\mathcal N_vb=\sqrt{s}\,b$ and $\mathcal N_TT=sT$.
Changing variables in the normalized phase/spatial integral gives
\begin{equation}\label{stress-short:lifts}
 \|\mathcal N_vb\|_{Y_2}^2=\pi R^2\|b\|_H^2,\qquad
 \|\mathcal N_TT\|_{Y_2}^2=\pi R^2\|T\|_{H_T}^2.
\end{equation}
Indeed the physical space Jacobian after azimuthal integration is
$\pi R^2s^{-5/2}ds\,d\eta$ and normalized phase becomes $dq/c$.
The physical divergence and Leray projection remain in the free inverse;
these formulas only identify its Hilbert spaces.

Fix $m_Z>3/2$, for example $m_Z=2$, and define
\[
 W=\left\langle\eta/\sqrt{1+s}\right\rangle^{m_Z},\qquad
 Z_1=L^2(s^{-1/2}W^2d\varsigma),\qquad
 Z=Z_1\oplus\rho_b^{-1/2}H_T.
\]
The norm on $Z$ is the sum norm
$\|(T_1,T_2)\|_Z=\|T_1\|_{Z_1}+\rho_b^{-1/2}\|T_2\|_{H_T}$.
Here
\[
 \rho_b=P_\nu\epsilon^{3/11}e^{m_0/\nu}
\]
is a fixed upper bound from the exact-background construction,
chosen so that
$\sup_\tau(\|r_b(\tau)\|_{11/4}+\|r_b(\tau)\|_\infty)\le\rho_b$.
It is not the radial coordinate. The stress inverse acts on
$T_1+T_2$; the second entry is not assigned a transverse weight.

Write $v=s^{-1/2}V$ and $v_b=s^{-1/2}r_b$. The train heat expansion,
transverse bounds, and compact-age inlet construction imply
\begin{equation}\label{stress-short:coefficients}
 \begin{aligned}
 \sup\sqrt{s}\,W|v|&\le P_\nu,&
 \sup_{s>S}\sqrt{s}\,W|v|&\le P_\nu(S^{-1/2}+\epsilon),\\
 \sup_{s<\sigma}\sqrt{s}\,W|v|&\le P_\nu\sqrt{\sigma},
       &&0<\sigma\le1,\\
 \|Wv(s)\|_{L^2_\eta L^\infty_q}
       &\le P_\nu[(1+s)^{-3/4}
                         +\epsilon(1+s)^{-1/4}].
 \end{aligned}
\end{equation}
Here $S$ is larger than the compact-age inlet support. To check the
age powers, the mean stream and its fixed weighted derivatives have
size $(1+s)^{-1/2}$ on transverse scale $\sqrt{1+s}$.
One transverse derivative gives the velocity size $(1+s)^{-1}$.
The stream completion
$\epsilon(1-\eta\partial_\eta-2s\partial_s)\psi$ has size
$P_\nu\epsilon(1+s)^{-1/2}$. Its slower decay is retained.
The $L^2_\eta$ norm adds the factor $(1+s)^{1/4}$.
Fluctuations have exponential decay after overlap; before overlap
the Gaussian comparison gives the same polynomial bounds.
The finite inlet is absorbed by $P_\nu$, and early smoothness gives
the third bound. Multiplying by any further fixed power of
$\langle\eta/\sqrt{1+s}\rangle$ leaves these conclusions unchanged,
by the weighted train expansion.

If $S=\epsilon^{-b}$ and $D_\perp>0$ is fixed, the coefficient can
also be localized to
\[
 |\eta|\le L_\perp=C_\perp(c/\nu+\sqrt S+1)
\]
with
\begin{equation}\label{stress-short:transverse}
 \sup_{s<S,\ |\eta|>L_\perp}
 \sqrt{s}\,W(|v|+|v_{\rm flat}|)
       \le e^{-D_\perp/\nu}.
\end{equation}
The Gaussian vorticity tail controls the heat-width part.
The nonzero horizontal stream kernels
$e^{-|k_n||\eta|}/(2|k_n|)$, $|k_n|\ge2\pi/c$, control the
$c/\nu$ part. The zero mode has decaying heat primitives.
Increasing $C_\perp$ absorbs the fixed weights, derivative budgets,
and the finite inlet's stream tail. No such localization is claimed
for $r_b$.

Define the exact split multiplier by
\[
 \mathcal M_Ub=(v\otimes b+b\otimes v,\
                         v_b\otimes b+b\otimes v_b).
\]
The flat multiplier uses the actual planar train in its first entry
and zero in its second entry. Since the ratio of stress and velocity
age densities is $s$, \eqref{stress-short:coefficients} gives
\begin{equation}\label{stress-short:multiplier}
 \|\mathcal M_Ub\|_Z\le P_\nu\|b\|_H,\qquad
 \|(\mathcal M_Ub)_2\|_{\rho_b^{-1/2}H_T}
                         \le2\sqrt{\rho_b}\|b\|_H.
\end{equation}
Its thin old and young restrictions have respectively norms
$P_\nu(S^{-1/2}+\epsilon)$ and $P_\nu\sqrt{\sigma}$.
Indeed $|a\otimes b+b\otimes a|\le2|a||b|$, so the thin
multiplier norm is at most $2\sup\sqrt{s}W|v|$.
For the repair, $\sqrt{s}v_b=r_b$ and multiplication followed by
the factor $\rho_b^{-1/2}$ gives the second estimate.
These bounds also hold after multiplication of both sides by any
common positive age weight, in particular the gain $g(s)^{-1}$.

For $A\ge1$ the thin source has the absolute old moment bound
\begin{equation}\label{stress-short:moment}
 \mathcal N_1(1_{s>A}T):=
 \int_A^\infty s^{-3/2}\int
 (1+|\eta|/\sqrt{s})\|T(s,\cdot,\eta)\|_{L^2(dq/c)}\,d\eta\,ds
 \le C_{m_Z}A^{-1/2}\|1_{s>A}T\|_{Z_1}.
\end{equation}
Transverse Cauchy--Schwarz gives
\[
 \int(1+|\eta|/\sqrt{s})\|T(s,\cdot,\eta)\|_{L^2_q}\,d\eta
       \le C_{m_Z}s^{1/4}\|WT(s)\|_{L^2_{q,\eta}},
\]
because $\int\langle x\rangle^{2-2m_Z}dx<\infty$.
The remaining time integral is
$\int_A^\infty s^{-5/4}\|WT(s)\|_2ds$.
Cauchy--Schwarz against $s^{-1/2}ds$ leaves
$\int_A^\infty s^{-2}ds=A^{-1}$, proving
\eqref{stress-short:moment}.

In outer coordinates $x'=y'/R$, $u=|x'|^{-2}$ and
$Y=\eta/\sqrt u$, the thin physical tensor is represented by
$F=uT$. Directly,
\[
 \int(1+|Y|)\|F(x',Y,\cdot)\|_{L^2(dq/c)}\,dx'\,dY
       =\pi\mathcal N_1(T).
\]
Thus \eqref{stress-short:moment} supplies exactly the first moment
used in the separated-kernel approximation. It also justifies
extending an old annulus to the axis by moment-norm approximation.
All cutoffs multiply the tensor before its distributional divergence;
no artificial boundary force is removed. The broad repair is instead
estimated by its physical $L^{11/4}$ coefficient and energy bounds.
